\section{模型假设}

\begin{enumerate}
  \item \textbf{源静止、位置固定未知。} 本问只讨论一个已测得示向度的干扰源，其在圆域内的
        位置 $G$ 在一局内不变；不引入源的移动或漂移。
  \item \textbf{测向误差按有界区间处理。} 示向度读数与真实方向的偏差落在 $[-1^\circ,1^\circ]$
        内，\textbf{不假设其概率分布}。这与题面"误差不超过 $1^\circ$"一致，也使结论成为
        保证型结论（对应文献中区间/集员估计一脉）。
  \item \textbf{接收半径固定未知。} 每个源的接收半径 $\rho_{\mathrm{rec}}\in[1000,1500]$ m
        在一局内固定但不可观测；"听到某源"等价于"到该源的距离不超过 $\rho_{\mathrm{rec}}$"。
        本问只研究能听到的情形（对定向源，认为机器狗位于迎光侧或源为全向源）。
  \item \textbf{由第一次测向推出的源距离区间。} $S_1$ 处测出了示向度，故源必满足
        $d=|S_1G|>5$ m（附件规定距离不超过 $5$ m 时信号过强、测不出示向度）且
        $d\le1500$ m（接收半径上限）；同时源位于 $1800$ m 圆域内。这三条构成本文的不确定集 $C$。
  \item \textbf{定位区域口径与问题一一致。} 两条示向度各带 $\pm1^\circ$ 误差，其边界射线
        围成的凸多边形（被 $1800$ m 圆域裁剪）即定位区域；定位精度用该区域的\textbf{直径}
        度量。这一口径与附录图 2 的"两条射线夹出的定位四边形"一致。
  \item \textbf{可测性取最保守口径。} 要求第二次测向\textbf{一定}成功，即对不确定集内任意
        源都有 $|S_2G|\le1000$ m（接收半径下界）。若允许"第二次可能失败、需要第三次"，
        则改用 $1500$ m，本文在灵敏度分析中给出该口径的结果对照（§\ref{sec:sens}）。
  \item \textbf{清除成功的充分条件。} 清除要求机器狗进入距源 $20$ m 以内。若定位区域直径
        $D\le20$ m，把机器狗开到该区域的最小覆盖圆心，由平面集合的 Jung 定理可知区域内
        任意点到该圆心距离不超过 $D/\sqrt3\le11.5$ m $<20$ m，故清除必然成功。本文因此把
        "$20$ m 硬阈值"落实为对定位区域直径的要求，并始终以"最坏直径尽可能小"为选点目标。
  \item \textbf{只考虑几何与信息层面的代价。} 选点结论与机器狗的移动路径、能耗、地形等无关
        （这些属于问题三、四的规划层面）；第二检测点选定后的实际执行路线不在本问范围内。
  \item \textbf{数值离散足够细。} 后文所有连续量都通过确定性网格离散近似，网格步长造成的
        误差远小于 $\pm1^\circ$ 测向误差对应的几十米量级，且通过加密 $3$ 倍复算校验
        （偏差 $6\times10^{-15}$ 相对量级，见 §\ref{sec:eval}）。
\end{enumerate}
